\begin{helper}
Fix finite source and target vertex sets, a source common coordinate set
\(S\subseteq V\), a target common coordinate set \(S'\subseteq V'\), and an
injective transport map \(\phi:V\to V'\) with \(S'=\phi(S)\).  Assume the
source and target activation--priority laws have the Bernoulli activation
atom formula with parameter \(\gamma/\Delta\in[0,1]\) and the finite priority
lower-rectangle formula.

Let
\[
  \Omega_0=\{A\subseteq V:A\subseteq S\}
\]
index the common activation trace.  For \(A\in\Omega_0\), set
\[
  B_A=\{\eta:\eta_{\rm act}\cap S=A\},\qquad
  B'_A=\{\eta':\eta'_{\rm act}\cap S'=\phi(A)\}.
\]
Let
\[
  \Lambda=\{0,1\}\times S.
\]
The two labels attached to \(v\in S\) are
\[
  0<\pi(v),\qquad \pi(v)\le 1,
\]
and on the target side they are transported to
\[
  0<\pi'(\phi(v)),\qquad \pi'(\phi(v))\le 1.
\]
Let \(T_A(\eta,\eta')\) mean that
\[
  \eta'_{\rm act}\cap S'=(\eta_{\rm act}\cap S)^\phi
  \quad\text{and}\quad
  \pi'(\phi(v))=\pi(v)\ \text{for all }v\in S.
\]
Then the source and target base cells and label events are measurable, the
base cells are pairwise disjoint on each side, \(T_A\) transports the paired
base cells and every paired label, and every truth-vector atom
\[
  B_A\cap\{\eta:\eta\in Q_\lambda\Longleftrightarrow \lambda\in\tau
        \text{ for all }\lambda\in\Lambda\}
\]
has the same nonnegative real mass as its transported target atom.  Moreover
the base cells cover the full unit-priority support on each side, and the
complement of unit-priority support has measure zero for both laws.
\end{helper}
\begin{proof}
Write \(p=\gamma/\Delta\).  The hypotheses give \(0\le p\le1\), the exact
Bernoulli mass of every full activation atom, and, conditional on such a full
activation atom, the exact mass of every lower rectangle
\(\{0\le\pi(v)\le t_v\}_{v\in V}\) with \(0\le t_v\le1\).  The same three
facts hold on the target side with \(V'\).
The single-law support and boundary calculation is registered as
\noderef{SamplingPriorityPositiveUnitSupport}; its finite-fibre hypothesis
holds here because every activation mass is an \(\operatorname{ofReal}\)
value. The activation-trace calculation is registered as
\noderef{SamplingBernoulliActivationTraceMass}. We give the resulting
set and mass arguments explicitly below, keeping all atoms raw rather than
redefining them by intersection with the support.

We first record the elementary measurable and set-theoretic parts.  The
activation coordinate has the discrete finite measurable space, so
\(\eta\mapsto\eta_{\rm act}\cap S\) and
\(\eta'\mapsto\eta'_{\rm act}\cap S'\) are measurable maps into finite
discrete spaces.  Hence each fibre \(B_A\) and \(B'_A\) is measurable.  For
each \(v\in S\), the priority coordinate maps
\(\eta\mapsto \pi_\eta(v)\) and
\(\eta'\mapsto\pi_{\eta'}(\phi(v))\) are measurable real coordinate maps.
Thus the inverse images of \((0,\infty)\) and \((-\infty,1]\) give
measurability of all source and target label events.

The source base cells are disjoint because an activation trace on \(S\) has a
unique value.  If a target sample lies in \(B'_A\cap B'_{A'}\), then
\(\phi(A)=\phi(A')\) as subsets of \(S'\).  Since \(\phi\) is injective, its
restriction to \(S\) is injective, so \(A=A'\); hence the target base cells
are disjoint.  If \(T_A(\eta,\eta')\) holds, then by definition
\[
  \eta'_{\rm act}\cap S'=(\eta_{\rm act}\cap S)^\phi
  \quad\text{and}\quad
  \pi_{\eta'}'(\phi(v))=\pi_\eta(v)\quad(v\in S).
\]
Therefore \(\eta\in B_A\) iff \(\eta'\in B'_A\), and for each label
\(\lambda=(b,v)\in\{0,1\}\times S\), the source truth condition attached to
\(\lambda\) is equivalent to the transported target truth condition.

We next prove the unit-support nullity, because the raw truth atoms are most
conveniently computed after removing the null set where some priority
coordinate is outside \([0,1]\).  Fix a full source activation atom
\(\{\eta_{\rm act}=F\}\).  Apply the priority lower-rectangle hypothesis with
the constant threshold \(t_v=1\) for every \(v\in V\).  The product of the
thresholds is \(1\), so
\[
 \nu\{\eta:\eta_{\rm act}=F,\ 0\le\pi_\eta(v)\le1\text{ for all }v\in V\}
 =\nu\{\eta:\eta_{\rm act}=F\}.
\]
The set on the left is contained in the activation atom on the right, hence
the difference
\[
 \{\eta:\eta_{\rm act}=F,\ \neg(0\le\pi_\eta(v)\le1
       \text{ for all }v\in V)\}
\]
has measure zero.  The finitely many full activation atoms partition the
activation coordinate, so finite additivity over their disjoint union gives
\(\nu(U^c)=0\), where
\[
  U=\{\eta:0\le\pi_\eta(v)\le1\text{ for all }v\in V\}.
\]
The identical argument with \(V'\), \(\nu'\), and constant threshold \(1\)
gives \(\nu'((U')^c)=0\) for the target unit-priority support.

The same lower-rectangle law also gives the boundary nullities needed for the
strict labels.  Fix \(v_0\in V\) and a full activation atom \(F\).  Put
\(t_{v_0}=0\) and \(t_v=1\) for \(v\ne v_0\).  All thresholds lie in
\([0,1]\), and the product is \(0\), so
\[
 \nu\{\eta:\eta_{\rm act}=F,\ 0\le\pi_\eta(v)\le t_v
       \text{ for all }v\in V\}=0.
\]
This lower rectangle contains, inside the full unit cube, the boundary
\(\{\pi_\eta(v_0)=0\}\).  Summing over \(F\) gives
\(\nu(U\cap\{\pi(v_0)=0\})=0\).  Repeating this for the finitely many
\(v_0\in S\), and similarly for \(\phi(v_0)\in S'\), shows that all source
and target coordinate-zero boundaries used by the labels have measure zero.

Now fix \(A\in\Omega_0\) and a truth vector
\(\tau\subseteq\Lambda\).  Let \(Q_{A,\tau}\) be the corresponding source
truth-vector atom inside \(B_A\), and \(Q'_{A,\tau}\) the transported target
atom inside \(B'_A\).  Since \(U^c\) and \((U')^c\) have measure zero, it is
enough to compute \(Q_{A,\tau}\cap U\) and \(Q'_{A,\tau}\cap U'\).

Consider one coordinate \(v\in S\).  There are four possible truth
assignments for the two labels \(0<\pi(v)\) and \(\pi(v)\le1\).
Inside \(U\), if both labels are true, the condition is
\[
  0<\pi(v)\le1,
\]
which differs from \(0\le\pi(v)\le1\) only by the already null boundary
\(\pi(v)=0\).  If \(0<\pi(v)\) is false and \(\pi(v)\le1\) is true, then
inside \(U\) we have \(\pi(v)=0\), so this branch has measure zero.  If
\(0<\pi(v)\) is true and \(\pi(v)\le1\) is false, then inside \(U\) it would
require \(\pi(v)>1\), impossible.  If both labels are false, it requires both
\(\pi(v)\le0\) and \(\pi(v)>1\), impossible even before intersecting with
\(U\).  Thus, after discarding null coordinate-zero boundaries, the only
nonzero source atom is the one in which for every \(v\in S\) both labels are
true; every other truth vector has source mass zero.  The same four-case
analysis applies to the target coordinate \(\phi(v)\), and injectivity of
\(\phi\) identifies the list of target coordinates with the list of source
coordinates, so the zero and nonzero cases are transported exactly.

It remains to compute the nonzero case.  Suppose \(\tau\) contains both
labels attached to every \(v\in S\).  Modulo the zero coordinate boundaries,
\[
  Q_{A,\tau}\cap U
  =
  \{\eta:\eta_{\rm act}\cap S=A,\ 0\le\pi_\eta(v)\le1
       \text{ for all }v\in V\}.
\]
This set is the disjoint finite union over full activation sets
\[
  F\subseteq V,\qquad F\cap S=A,
\]
of the full activation atom \(\{\eta_{\rm act}=F\}\) intersected with the
full priority cube.  For each such \(F\), the constant-\(1\) rectangle
calculation above says that the priority cube contributes the full
activation mass, namely
\[
  \operatorname{ofReal}\bigl(p^{|F|}(1-p)^{|V|-|F|}\bigr).
\]
Since \(0\le p\le1\), these summands are nonnegative real masses.  Summing
over all \(F\) with trace \(A\) factors as
\[
  p^{|A|}(1-p)^{|S|-|A|}
  \sum_{E\subseteq V\setminus S} p^{|E|}(1-p)^{|V\setminus S|-|E|}.
\]
The finite binomial expansion of \((p+(1-p))^{|V\setminus S|}\), together
with \(p+(1-p)=1\), makes the outside sum equal to \(1\).  Hence the source
mass is represented by the nonnegative real
\[
  w_A=p^{|A|}(1-p)^{|S|-|A|}.
\]

On the target side, the nonzero transported atom is the disjoint finite union
over all \(F'\subseteq V'\) with
\[
  F'\cap S'=\phi(A).
\]
The same calculation gives
\[
  p^{|\phi(A)|}(1-p)^{|S'|-|\phi(A)|}.
\]
Because \(S'=\phi(S)\) and \(\phi\) is injective on \(S\), we have
\(|S'|=|S|\) and \(|\phi(A)|=|A|\).  Therefore the target real mass is the
same \(w_A\).  Thus in the nonzero truth-vector case we take
\(w=w_A\), while in every other truth-vector case we take \(w=0\).  The
preceding paragraph proves \(0\le w\) and
\[
  \operatorname{ofReal}(w)=\nu(Q_{A,\tau})
  \quad\text{and}\quad
  \operatorname{ofReal}(w)=\nu'(Q'_{A,\tau}),
\]
including the inconsistent-label and boundary cases.

Finally, the support-cover assertions are immediate from the definition of
the base cells.  If \(\eta\in U\), then
\(A=\eta_{\rm act}\cap S\) is a subset of \(S\), so \(A\in\Omega_0\) and
\(\eta\in B_A\).  If \(\eta'\in U'\), then
\(A'=\eta'_{\rm act}\cap S'\) is a subset of \(S'=\phi(S)\).  Since
\(\phi:S\to S'\) is a bijection onto its image, there is a unique
\(A\subseteq S\) with \(\phi(A)=A'\), and \(\eta'\in B'_A\).  Together with
the already proved \(\nu(U^c)=0\) and \(\nu'((U')^c)=0\), this establishes
all asserted measurability, disjointness, transport, truth-atom mass,
support-cover, and support-null conclusions.
\end{proof}
